\documentclass[aspectratio=169]{beamer}

\usepackage[utf8]{inputenc}
\usepackage[T1]{fontenc}
\usepackage[french]{babel}

\usetheme{Madrid}

% ── Palette bordeaux / rouge ─────────────────────────────────────────────────
\definecolor{bordeaux}{RGB}{112, 18, 28}
\definecolor{rouge}{RGB}{170, 40, 40}
\definecolor{beigelight}{RGB}{248, 238, 230}

\setbeamercolor{palette primary}{bg=bordeaux, fg=white}
\setbeamercolor{palette secondary}{bg=rouge, fg=white}
\setbeamercolor{palette tertiary}{bg=bordeaux!85!black, fg=white}
\setbeamercolor{palette quaternary}{bg=bordeaux, fg=white}
\setbeamercolor{structure}{fg=bordeaux}
\setbeamercolor{frametitle}{bg=bordeaux, fg=white}
\setbeamercolor{title}{bg=bordeaux, fg=white}
\setbeamercolor{section in toc}{fg=bordeaux}
\setbeamercolor{item}{fg=bordeaux}

\setbeamercolor{block title}{bg=bordeaux, fg=white}
\setbeamercolor{block body}{bg=beigelight, fg=black}
\setbeamercolor{block title alerted}{bg=rouge, fg=white}
\setbeamercolor{block body alerted}{bg=rouge!12, fg=black}
\setbeamercolor{block title example}{bg=bordeaux!70!black, fg=white}
\setbeamercolor{block body example}{bg=bordeaux!10, fg=black}

% ── Fondu sur les éléments non encore révélés ────────────────────────────────
\setbeamercovered{transparent=15}

\setbeamertemplate{navigation symbols}{}

\title[Winevizer]{Gestion de la relation client et rétention}
\subtitle{Cas n\textsuperscript{o}7 -- Winevizer}
\author{Groupe X}
\date{Juin 2026}
\institute{MGEHD2142 -- Management de la PME}

\begin{document}

% ============================================================
% PAGE DE GARDE
% ============================================================
\begin{frame}[plain,noframenumbering]
  \transfade
  \titlepage
\end{frame}

% ============================================================
% SLIDE 1 — INTRODUCTION
% ============================================================
\begin{frame}{Introduction}
  \transfade
  \begin{columns}
    \begin{column}{0.58\textwidth}
      \textbf{Winevizer} : sommelier virtuel SaaS B2B pour l'Horeca\\
      {\small QR code ou tablette, 7 langues, 24,99\,€/mois, fondée en 2021}
      \vspace{0.6em}

      \uncover<2->{
      Demande initiale : tunnel de vente, échecs d'abonnement, fidélisation.\\[0.4em]
      \textbf{Notre recentrage :} le churn des payants reste contenu ($\sim$10\,\% sur 3~ans).\\
      Le vrai frein est \textbf{l'activation}.
      }
    \end{column}
    \begin{column}{0.40\textwidth}
      \uncover<3->{
      \begin{block}{Chiffres clés}
        \begin{itemize}
          \item 585 comptes créés
          \item 62 souscriptions payantes
          \item \textbf{Conversion globale $\approx$ 10,6\,\%}
          \item Conversion fin d'essai (carte déployée) : \textbf{20--25\,\%}
        \end{itemize}
      \end{block}
      }
    \end{column}
  \end{columns}
\end{frame}

% ============================================================
% SLIDE 2 — PARTIE 2 : VPC — PROFIL CLIENT
% ============================================================
\begin{frame}{Partie 2 -- Value Proposition Canvas : Profil client}
  \transfade
  {\small \textbf{Le restaurateur indépendant / patron de bar à vins}}
  \vspace{0.2em}
  \begin{columns}
    \begin{column}{0.33\textwidth}
      \uncover<1->{
      \begin{block}{Customer Jobs}
        \scriptsize
        \begin{itemize}
          \item Maximiser la marge vins\\(30--40\,\% du CA)
          \item Service homogène malgré la rotation du personnel
          \item Carte à jour en temps réel
          \item Image professionnelle
        \end{itemize}
      \end{block}
      }
    \end{column}
    \begin{column}{0.33\textwidth}
      \uncover<2->{
      \begin{alertblock}{Pains}
        \scriptsize
        \begin{itemize}
          \item Pénurie de sommeliers\\(\textless{}5\,\% des établissements)
          \item Personnel peu formé\\(étudiants, flexi-jobs)
          \item Réimpression coûteuse des cartes
          \item Conseil multilingue impossible
        \end{itemize}
      \end{alertblock}
      }
    \end{column}
    \begin{column}{0.33\textwidth}
      \uncover<3->{
      \begin{exampleblock}{Gains attendus}
        \scriptsize
        \begin{itemize}
          \item Ticket moyen plus élevé
          \item Moins de dépendance à la qualification
          \item Meilleure satisfaction client
          \item Pilotage de la cave par la donnée
        \end{itemize}
      \end{exampleblock}
      }
    \end{column}
  \end{columns}
\end{frame}

% ============================================================
% SLIDE 4 — PARTIE 2 : VPC — FIT ET MISFIT
% ============================================================
\begin{frame}{Partie 2 -- Value Proposition Canvas : Fit et Misfit}
  \transfade
  \begin{columns}
    \begin{column}{0.5\textwidth}
      \uncover<1->{
      \begin{exampleblock}{Fit stratégique validé}
        \small
        Winevizer répond directement aux pains et jobs :
        \begin{itemize}
          \item Se substitue au sommelier humain
          \item Carte synchronisée, multilingue
          \item Pilotage par la donnée
        \end{itemize}
        \vspace{0.3em}
        Utilisateurs ayant \textbf{déployé} la carte : \textbf{20--25\,\% convertissent}.
      \end{exampleblock}
      }
    \end{column}
    \begin{column}{0.5\textwidth}
      \uncover<2->{
      \begin{alertblock}{Misfit : la valeur est accessible trop tard}
        \small
        Le problème n'est pas \textit{ce que} fait Winevizer, mais \textit{quand} et \textit{comment} la valeur devient accessible.
        \begin{itemize}
          \item Discours centré fonctionnalités, pas sur le ROI
          \item Signaux de confiance peu visibles
          \item Friction cumulative avant la 1\textsuperscript{re} valeur perçue
        \end{itemize}
      \end{alertblock}
      }
    \end{column}
  \end{columns}
\end{frame}

% ============================================================
% SLIDE 5 — PARTIE 2 : PORTER — FRICTIONS DU PARCOURS CLASSIQUE
% ============================================================
\begin{frame}{Partie 2 -- Chaîne de valeur : Parcours classique}
  \transfade
  \textbf{6 étapes cumulatives avant la première valeur perçue}
  \vspace{0.15em}
  \begin{enumerate}
    \small
    \item<1-> \textbf{Création de compte :} 5 champs ; règles du mot de passe affichées \textit{après} saisie
    \item<2-> \textbf{Confirmation email :} bug de lien expiré ; étape contournable sans bénéfice
    \item<3-> \textbf{Création de l'établissement :} 8 champs obligatoires sans lien avec les fonctionnalités
    \item<4-> \textbf{Encodage de la carte :} import PDF -- méthode la plus naturelle -- \textbf{verrouillé derrière le paywall}
    \item<5-> \textbf{Traitement IA :} délai de plusieurs heures, \textbf{aucune valeur intermédiaire délivrée}
    \item<6-> \textbf{Génération du QR code :} non guidée, difficile à localiser dans l'interface
  \end{enumerate}
  \vspace{0.2em}
  \uncover<6->{\footnotesize \textit{S'ajoutent : défauts visuels sur la landing page ; chatbot présenté comme temps réel mais réponse asynchrone.}}
\end{frame}

% ============================================================
% SLIDE 6 — PARTIE 2 : PORTER — PARCOURS ALTERNATIF
% ============================================================
\begin{frame}{Partie 2 -- Chaîne de valeur : Parcours alternatif}
  \transfade
  {\small \textbf{Via la landing page :} email + nom d'établissement + fichier $\rightarrow$ carte générée en $\sim$2\,h}
  \vspace{0.4em}
  \begin{columns}
    \begin{column}{0.5\textwidth}
      \uncover<1->{
      \begin{exampleblock}{Avantages}
        \begin{itemize}
          \item Beaucoup moins de friction
          \item Import PDF librement accessible
          \item Aligné sur le besoin réel du restaurateur
        \end{itemize}
      \end{exampleblock}
      }
    \end{column}
    \begin{column}{0.5\textwidth}
      \uncover<2->{
      \begin{alertblock}{Incohérences et fragilités}
        \begin{itemize}
          \item PDF \textbf{gratuit ici}, payant dans le parcours classique
          \item Lien email = unique porte d'accès : suppression accidentelle $\Rightarrow$ import perdu
          \item Aucune orientation vers ce parcours depuis le site
        \end{itemize}
      \end{alertblock}
      }
    \end{column}
  \end{columns}
\end{frame}

% ============================================================
% SLIDE 7 — PARTIE 3 : RECOMMANDATIONS — LANDING PAGE
% ============================================================
\begin{frame}{Partie 3 -- Recommandations : Landing page (R1--R6)}
  \transfade
  \begin{columns}
    \begin{column}{0.5\textwidth}
      \uncover<1->{
      \textbf{R1 -- Discours centré sur le ROI}\\
      \small Afficher les indicateurs d'impact dès la page d'accueil, comme les concurrents (Somm'it : \og{}+3\,pts de marge\fg{} ; Coena : \og{}+20\,\% de CA vins\fg{})
      \vspace{0.5em}

      \textbf{R2 -- Preuve sociale crédible}\\
      \small Témoignages nominatifs avec logos, remontés en section principale
      \vspace{0.5em}

      \textbf{R3 -- Afficher les intégrations POS}\\
      \small Logos des POS compatibles visibles dès la page d'accueil
      }
    \end{column}
    \begin{column}{0.5\textwidth}
      \uncover<2->{
      \textbf{R4 -- Corriger les défauts visuels}\\
      \small Animations saccadées, vidéo débordante $\rightarrow$ nuit à l'impression de qualité
      \vspace{0.5em}

      \textbf{R5 -- Supprimer les sous-pages redondantes}\\
      \small Rapatrier ROI et témoignages directement en page principale
      \vspace{0.5em}

      \textbf{R6 -- Simplifier la présentation tarifaire}\\
      \small Afficher explicitement l'économie du plan annuel : \og{}249,90\,€, soit 49,98\,€ économisés\fg{}
      }
    \end{column}
  \end{columns}
\end{frame}

% ============================================================
% SLIDE 8 — PARTIE 3 : RECOMMANDATIONS — ONBOARDING
% ============================================================
\begin{frame}{Partie 3 -- Recommandations : Onboarding (R7--R12)}
  \transfade
  \begin{columns}
    \begin{column}{0.5\textwidth}
      \uncover<1->{
      \textbf{R7 -- Alléger le formulaire d'inscription}\\
      \small Supprimer Prénom/Nom ; afficher/masquer le MDP ; règles avant la saisie
      \vspace{0.5em}

      \textbf{R8 -- Magic link / OTP sans mot de passe}\\
      \small Connexion sans mot de passe (comme Slack, Notion) ; corriger le bug du lien expiré
      \vspace{0.5em}

      \textbf{R9 -- Persistance de session}\\
      \small Cookie sécurisé (7--30 jours) ; évite les reconnexions systématiques en service
      }
    \end{column}
    \begin{column}{0.5\textwidth}
      \uncover<2->{
      \textbf{R10 -- Alléger la création d'établissement}\\
      \small Réduire à 2 champs essentiels (nom + type) ; le reste en optionnel
      \vspace{0.5em}

      \textbf{R11 -- Lever le paywall sur l'import PDF}\\
      \small Accessible dès la période d'essai ; le fondateur envisage déjà cette décision
      \vspace{0.5em}

      \textbf{R12 -- Rationaliser / unifier les parcours}\\
      \small Faire du parcours alternatif le point d'entrée principal ; corriger la fragilité du lien email unique
      }
    \end{column}
  \end{columns}
\end{frame}

% ============================================================
% SLIDE 9 (= dernière) — PARTIE 3 : RECOMMANDATIONS TECHNIQUE & CONCLUSION
% ============================================================
\begin{frame}{Partie 3 -- Recommandations techniques \& Conclusion}
  \transfade
  \begin{columns}
    \begin{column}{0.46\textwidth}
      \uncover<1->{
      \begin{block}{Recommandations techniques (R13--R15)}
        \small
        \begin{itemize}
          \item \textbf{R13} Automatiser le chatbot (FAQ, QR code, imports)
          \item \textbf{R14} Corriger le bug CSV (1\textsuperscript{re} ligne ignorée)
          \item \textbf{R15} Bloquer la réactivation de l'essai gratuit\\(email chiffré conservé en base)
        \end{itemize}
      \end{block}
      }
    \end{column}
    \begin{column}{0.52\textwidth}
      \uncover<2->{
      \begin{block}{Ce que montre l'analyse}
        \small
        \begin{itemize}
          \item Fit produit-marché \textbf{solide} : 20--25\,\% convertissent quand la valeur est atteinte
          \item Le frein tient à trois leviers : \textbf{discours}, \textbf{frictions} et \textbf{incohérences}
        \end{itemize}
      \end{block}
      \vspace{0.4em}
      \begin{alertblock}{Priorités}
        \small
        \textbf{Court terme :} R1--R6, R11, R14\\
        \textbf{Moyen terme :} R7--R12\\
        \textit{Hors périmètre : échecs d'abonnement, fidélisation des payants}
      \end{alertblock}
      }
    \end{column}
  \end{columns}
\end{frame}

\end{document}
